% !PRINTED 649
\weberpage{649}
As a foundation we here use the function
\weq{1}{w = (k_0 - k_2)(k_1 - k_3),}
which must be different from zero, and which for a real field $\mathfrak{K}$
is always real, since either $k_0, k_1, k_2, k_3$ are real or $k_0, k_2$ and
$k_1, k_3$ form two conjugate imaginary pairs.

Under the cyclic permutation $(0, 1, 2, 3)$, $w$ changes its sign. Hence the
square of $w$ is a cyclic function, which by hypothesis is to be known. We
therefore set, denoting by the small latin letters $a, b, c \ldots$
quantities in $\mathfrak{K}$, that is, rational quantities,
\weq{2}{w = 4\sqrt{c},}
and remark further that every function of the $k$ which changes its sign
under the cyclic permutation $(0, 1, 2, 3)$ is the product of a rational
quantity with $\sqrt{c}$.

It is then only a matter of forming cyclic functions in sufficient number, so
that the four quantities $k_0, k_1, k_2, k_3$ can be computed from them and
expressed algebraically in terms of mutually independent rational quantities.

Now two further cyclic functions are
\[
(k_0 - k_2)^2 + (k_1 - k_3)^2, \qquad \frac{(k_0 - k_2)^2 - (k_1 - k_3)^2}{(k_0 - k_2)(k_1 - k_3)},
\]
and we thus obtain, if $a, b$ are rational quantities,
\weq{3}{\begin{aligned}
(k_0 - k_2)^2 + (k_1 - k_3)^2 &= 8\,b \\
(k_0 - k_2)^2 - (k_1 - k_3)^2 &= 8\,a\sqrt{c}.
\end{aligned}}
The three rational quantities $a, b, c$ are, however, not independent of one
another; rather, by (2) and (3) there subsists between them the relation
\weq{4}{b^2 = c\,(1 + a^2).}
From (3) it follows further that
\weq{5}{\begin{aligned}
k_0 - k_2 &= 2\sqrt{b + a\sqrt{c}} \\
k_1 - k_3 &= 2\sqrt{b - a\sqrt{c}},
\end{aligned}}
and between the two square roots there subsists the relation
\weq{6}{\sqrt{c} = \sqrt{b + a\sqrt{c}}\; \sqrt{b - a\sqrt{c}}.}
If we further denote by $C$ and $B$ again two rational quantities, then
%CONT>>
